% Part: proof-theory
% Chapter: proof-search
% Section: completeness

\documentclass[../../../include/open-logic-section]{subfiles}

\begin{document}

\olfileid[pt]{pt}{ps}{cpl}

\olsection{Completude}

Mostramos agora que o algoritmo de busca de provas é infalível, isto é,
que, se ele abortar ou nunca parar, não haverá !!{proof} do sequente
inicial~$\Gamma \Sequent \Delta$. Para isso, definimos
!!a{structure}~$\Struct M$ tal que, para todo $!A \in \Gamma$,
$\Sat{M}{!A}$, e, para todo~$! \in \Delta$, $\Sat/{M}{!A}$. Em outras
palavras, todas as !!{formula}s de~$\Gamma$ são verdadeiras e todas as
!!{formula}s de~$\Delta$ são falsas em~$\Struct M$, isto é, $\Gamma
\Sequent \Delta$ não é válido. Como \Log{G3c} é correto, $\Gamma
\Sequent \Delta$ não tem !!{proof}

Definimos $\Struct{M}$ a partir de um par $\Theta, \Xi$ de
!!{formula}s e de um conjunto~$C$ de !!{constant}s. Obtemos $\Theta,
\Xi$ e~$C$ de um \emph{ramo de falha} da execução malsucedida (abortada
ou não terminante) da busca de provas. Um ramo de falha é uma sequência
$\Pi_n \Sequent \Lambda_n$ de sequentes e conjuntos de
!!{constant}s~$C_n$ tal que $\Pi_0 \Sequent \Lambda_0$ é o sequente
final $\Gamma \Sequent \Delta$, e $\Pi_{n+1} \Sequent \Lambda_{n+1}$
é uma premissa de $\Pi_n \Sequent \Lambda_n$ gerada no estágio~$n+1$
para a qual o algoritmo não encontra !!a{proof}. É evidente que uma
premissa desse tipo deve existir para cada~$n$, pois, do contrário, o
algoritmo teria encontrado !!a{proof}. Seja $C_n$ o conjunto de
!!{constant}s associado a~$\Pi_n \Sequent \Delta_n$.

Se o algoritmo abortar em um sequente superior constituído apenas por
!!{formula}s atômicas, o ramo que termina nesse sequente será um ramo
de falha. Se se o algoritmo nunca terminar, ele continuará gerando
árvores cada vez maiores. Podemos considerá-las o crescimento de uma
árvore infinita (a árvore que teríamos depois que infinitos passos do
algoritmo fossem executados). Essa seria uma árvore infinita, mas
finitamente ramificada (cada nó tem apenas um ou dois filhos---os
sequentes imediatamente acima dele). Pelo Lema de K\H{o}nig, toda
árvore infinita e finitamente ramificada deve ter um ramo infinito. Um
ramo infinito desse tipo seria um ramo de falha no caso em que o
algoritmo de busca fosse executado para sempre.

Se a sequência de $\Pi_n \Sequent \Lambda_n$ e $C_n$ for um ramo de
falha, sejam $\Theta = \bigcup_n \Pi_n$ e $\Xi = \bigcup_n \Lambda_n$
(removendo os índices e as ocorrências múltiplas de !!{formula}s), e
$C = \bigcup_n C_n$.

Estamos prontos para definir~$\Struct{M}$.
\begin{enumerate}
\item O domínio $\Domain{M}$ é o conjunto dos termos que podem ser
formados a partir de~$C$ e das !!{function}s em~$\Gamma \Sequent
\Delta$, se houver alguma, mas sem !!{variable}s.
\item Se $c \in \Domain{M}$, $\Assign{c}{M} = c$.
\item Se $f$ for uma !!{function} de $m$ lugares em $\Gamma \Sequent
\Delta$, e $t_1$, \dots, $t_m \in \Domain{M}$, então
$\Assign{f}{M}(t_1, \dots, t_m) = f(t_1, \dots, t_n)$.
\item Se $R$ for um !!{predicate} de $m$ lugares em $\Gamma \Sequent
\Delta$, então $\Assign{R}{M} = \Setabs{\tuple{t_1, \dots, t_m} \in
\Domain{M}^n}{R(t_1, \dots, t_n) \in \Theta}$.
\end{enumerate}
$\Struct{M}$ é um \emph{modelo de termos}, pois seu domínio é um
conjunto de termos, e as !!{constant}s e !!{function}s são interpretadas
de modo a garantir que $\Value{t}{M} = t$. As !!{formula}s atômicas
em~$\Theta$ são usadas para definir $\Assign{R}{M}$ de modo a garantir
que $\Sat{M}{R(t_1, \dots, t_n)}$ sse $R(t_1, \dots, t_n) \in
\Theta$.

\begin{lem}\ollabel{lem:truth}
Para todo $!A \in \Theta \cup \Xi$:
\begin{enumerate}
  \item Se $!A \in \Theta$, então $\Sat{M}{!A}$.
  \item Se $!A \in \Xi$, então $\Sat/{M}{!A}$.
\end{enumerate}
\end{lem}

\begin{proof}
  Por indução sobre $\depth{!A}$.

Primeiro, suponha que $!A$ seja atômica, isto é, que $!A \ident
\lfalse$ ou $!A \ident R(t_1, \dots, t_m)$.

Como $\Pi_n \Sequent \Lambda_n$ é um ramo de falha, nenhum $\Pi_n$
pode conter $\lfalse$ (caso contrário, $\Pi_n \Sequent \Lambda_n$
seria um axioma). $\Sat{M}{R(t_1, \dots, t_m)}$ sse $R(t_1, \dots,
t_m) \in \Theta$, pela definição de~$\Struct{M}$. Em conjunto, esses
fatos mostram que, se $!A \in \Theta$, então $\Sat{M}{!A}$.

Se $!A$ for atômica e $!A \in \Pi_n$, então $!A \in \Pi_{n+1}$; e,
se $!A \in \Lambda_n$, então $!A \in \Lambda_{n+1}$. (Isso decorre
da maneira como o algoritmo de busca de provas gera os sequentes
sucessores.) Portanto, se $!A \in \Theta \cap \Xi$, então, para algum
~$n$, $!A \in \Pi_n \cap \Lambda_n$. Isso significa que $\Pi_n
\Sequent \Lambda_n$ é um axioma, e o ramo de falha não pode conter
axiomas. Assim, por construção, $!A \notin \Theta \cap \Xi$ para
$!A$ atômica e, consequentemente, se $!A \in \Xi$, então $!A \notin
\Theta$. Isso implica que, se $!A \in \Xi$, então $\Sat/{M}{!A}$,
pela definição de~$\Struct{M}$.

Para o passo indutivo, suponha que (1) e (2) valham para todas as
!!{formula}s de !!{depth} menor que a de~$!A$. Provamos (1) e (2) para
$!A$ distinguindo casos.

Observe primeiro que, quando $!A^i$ ocorre em $\Pi_n \Sequent
\Lambda_n$, $!A^i$ também ocorre em $\Pi_{n+1} \Sequent
\Lambda_{n+1}$, a menos que $i$ seja o menor índice em $\Pi_n \Sequent
\Lambda_n$. Como o menor índice aumenta nos sequentes superiores
acrescentados em cada estágio, acabamos chegando a $\Pi_n \Sequent
\Lambda_n$ no qual $!A^i$ ocorre e $i$ é o menor índice. No estágio
$n+1$, o algoritmo reduz~$!A^i$. Em outras palavras, se $!A \in
\Theta$, então, para algum $n$, $\Pi_n = \Pi_n', !A^i$ e $i$ é o
menor índice. E, se $!A \in \Xi$, então, para algum $n$, $\Lambda_n =
\Lambda_n', !A^i$ e $i$~é o menor índice.

\begin{enumerate}
  \item $!A \in \Theta$ e $!A \ident !B \land !C$. Então, para algum
  $n$, $\Pi_n = \Pi_n', (!B \land !C)^i$. $\Pi_{n+1} \Sequent
  \Lambda_{n+1}$ é a premissa correspondente de~\LeftR{\land}, isto é,
  $\Pi_{n+1} = \Pi_n', !B^k, !C^{k+1}$. Consequentemente, $!B, !C \in
  \Theta$. Pela hipótese de indução, $\Sat{M}{!B}$ e $\Sat{M}{!C}$.
  Pela definição de $\Sat{M}{}$, temos $\Sat{M}{!B \land !C}$.

  \item $!A \in \Xi$ e $!A \ident !B \land !C$. Então, para algum $n$,
  $\Lambda_n = \Lambda_n', !B \land !C$. $\Pi_{n+1} \Sequent
  \Lambda_{n+1}$ é uma das duas premissas de~\RightR{\land} com
  conclusão $\Pi_n \Sequent \Lambda'_n, !B \land !C$, isto é,
  $\Lambda_{n+1} = \Lambda_n', !B^k$ ou $\Lambda_{n+1} = \Lambda_n',
  !C^k$. Consequentemente, ou $!B \in \Xi$, ou $!C \in \Xi$. Pela
  hipótese de indução, ou $\Sat/{M}{!B}$, ou $\Sat/{M}{!C}$. Pela
  definição de $\Sat{M}{}$, temos $\Sat/{M}{!B \land !C}$.
  
  \item Os casos em que $!A \ident \lnot !B$, $!A \ident !B \lor !C$
  e $!A \ident !B \lif !C$ são semelhantes e ficam como exercícios.

  \item $!A \in \Theta$ e $!A \ident \lexists[x][!B(x)]$. Então, para
  algum $n$, $\Pi_n = \Pi_n', \lexists[x][!B(x)]^i$. $\Pi_{n+1}
  \Sequent \Lambda_{n+1}$ é a premissa correspondente
  de~\LeftR{\lexists}, isto é, $\Pi_{n+1} = \Pi_n', !B(c)^k$ para
  algum~$c$. Consequentemente, $!B(c) \in \Theta$ e~$c \in C$. Pela
  hipótese de indução, $\Sat{M}{!B(c)}$. Pela definição de
  $\Sat{M}{}$, temos $\Sat{M}{\lexists[x][!B(x)]}$.

  \item $!A \in \Xi$ e $!A \ident \lexists[x][!B(x)]$. Então, para algum
  $n$, $\Lambda_n = \Lambda_n', \lexists[x][!B(x)]^i$. Como, cada vez
  que $\lexists[x][!B(x)]$ é reduzida, $\lexists[x][!B(x)]$ permanece no
  lado direito do sequente-premissa com um novo índice,
  $\lexists[x][!B(x)]$ é reduzida infinitas vezes no lado direito de
  um ramo de falha se ocorrer ali. Todo termo $t \in \Domain{M}$ acaba
  sendo ``o primeiro termo que contém apenas constantes em~$C_n$ ainda
  não usadas em uma redução de $\lexists[x][!B(x)]$ à direita'' em um
  ramo de falha. Portanto, para todo $t \in \Domain{M}$, $!B(t) \in
  \Lambda_n$ para algum~$n$ e, consequentemente, $!B(t) \in \Xi$.
  Pela hipótese de indução, $\Sat/{M}{!B(t)}$ para todo $t \in
  \Domain{M}$. Pela definição de $\Sat{M}{}$, temos
  $\Sat/{M}{\lexists[x][!B(x)]}$.

  \item Os casos em que $!A \ident \lforall[x][!B(x)]$ são semelhantes e
  ficam como exercícios.
\end{enumerate}
\end{proof}

\begin{cor}\ollabel{cor:complete} O algoritmo de busca de provas para
\Log{G3c} é completo: se abortar ou nunca parar, o sequente
inicial~$\Gamma \Sequent \Delta$ não será válido e, portanto, não terá
!!{proof}.
\end{cor}

\begin{proof}
  Se o algoritmo abortar ou nunca parar, haverá um ramo de falha a
  partir do qual $\Struct{M}$ poderá ser definida. Por \olref{lem:truth},
  para todo $!A \in \Gamma$, $\Sat{M}{!A}$ e, para todo $!A \in \Delta$,
  $\Sat/{M}{!A}$. (Lembre-se de que $\Pi_0 = \Gamma$ e $\Lambda_0 =
  \Delta$, logo $\Gamma \subseteq \Theta$ e $\Delta \subseteq \Xi$.)
  Portanto, $\Gamma \Sequent \Delta$ não será válido. Pela correção
  de~\Log{G3c}, $\Gamma \Sequent \Delta$ não terá !!{proof}.
\end{proof}

\begin{cor}
\Log{G3c} é completo: se $\Gamma \Sequent \Delta$ for válido,
$\Log{G3c} \Proves \Gamma \Sequent \Delta$.
\end{cor}

\begin{proof}
  Demonstramos a contrapositiva. Suponha $\Log{G3c} \Proves/ \Gamma
  \Sequent \Delta$. Então o algoritmo de busca de provas deve abortar
  ou nunca parar, pois não há !!{proof} a ser encontrada. Por
  \olref{cor:complete}, $\Gamma \Sequent \Delta$ não é válido.
\end{proof}


\end{document}
